\documentclass[11pt,a4paper]{article}
\usepackage[T1]{fontenc}
\usepackage{lmodern,amsmath,amssymb,amsthm,microtype}
\usepackage[margin=27mm]{geometry}
\usepackage[colorlinks=true,linkcolor=blue,citecolor=blue,urlcolor=blue]{hyperref}
\newtheorem{theorem}{Theorem}[section]
\newtheorem{proposition}[theorem]{Proposition}
\newtheorem{corollary}[theorem]{Corollary}
\newcommand{\Z}{\mathbb Z}\newcommand{\Q}{\mathbb Q}
\newcommand{\OK}{\mathcal O_K}
\setlength{\emergencystretch}{2em}
\title{Finite exceptions in Diophantine integrality criteria}
\author{}\date{17 November 2026}
\begin{document}\maketitle
\begin{abstract}
A Diophantine subset of a number field suffices to define its ring
of integers if it contains all but finitely many rational integers
and its nonintegral elements have bounded absolute height. The
criterion permits nonrational integral inputs. A finite correction
uses one additional coordinate, after which an integral basis gives
a definition of the whole ring of integers. A stronger hypothesis
that bounds every noninteger input defines the rational integers
as well. The result isolates the amount of height soundness needed
for an existence theorem; it does not construct the required
Diophantine subset.
\end{abstract}

\section{Finite corrections over a number field}
All Diophantine definitions in a number field $K$ allow
coefficients in $K$ and quantify every witness over $K$.
Every coordinate count includes the distinguished input.

\begin{proposition}\label{prop:correction}
Suppose $S\subseteq K$ is defined by
\[
 x\in S\quad\Longleftrightarrow\quad
 \exists\mathbf y\in K^{N-1}: P(x,\mathbf y)=0,
\]
where $P\in K[X,\mathbf Y]$ has total degree at most $D$.
Let $A,B\subseteq K$ be finite and disjoint. Then
$(S\setminus A)\cup B$ has a single-polynomial definition with
$N+1$ total coordinates and degree at most
\[
 |B|+2\max(D,|A|+1).
\]
\end{proposition}
\begin{proof}
Choose $\delta\in K$ that is not a square and put
\[
 a(X)=\prod_{\alpha\in A}(X-\alpha),\qquad
 b(X)=\prod_{\beta\in B}(X-\beta),
\]
with empty products equal to one. In the coordinates
$(X,\mathbf Y,V)$ use
\[
 b(X)\left(P(X,\mathbf Y)^2-
                 \delta\bigl(Va(X)-1\bigr)^2\right)=0.
\]
If $b(X)\ne0$, the nonsquare property of $\delta$ forces both
$P(X,\mathbf Y)=0$ and $Va(X)=1$. These conditions hold exactly
for $X\in S\setminus A$. If $b(X)=0$, the equation is satisfied
and $X\in B$. Conversely each of the two sets supplies witnesses.
The coordinate and degree bounds follow directly.
\end{proof}

No effective identification of $A$ or $B$ is implicit here.
Once the finite sets and the original polynomial are supplied,
the new polynomial is explicit.

\section{Bounded height of nonintegral inputs}
Let $h$ be the absolute logarithmic Weil height on
$\mathbb P^1(\overline{\Q})$. Northcott's theorem states that
algebraic numbers of bounded degree and bounded height form a
finite set~\cite[Theorem B.2.3]{HS}. In particular, a bounded-height
subset of a fixed number field is finite.

\begin{theorem}\label{thm:integers}
Let $K$ be a number field and let $S\subseteq K$ be Diophantine
in $K$. Suppose
\begin{enumerate}
\item $\Z\setminus S$ is finite;
\item there is a real constant $C$ such that every
      $x\in S\setminus\OK$ satisfies $h(x)\le C$.
\end{enumerate}
Then $\OK$ is Diophantine in $K$.
More precisely, set
\[
 A=S\setminus\OK,\qquad B=\Z\setminus S.
\]
These sets are finite and disjoint. If $S$ is given by $P$ as in
Proposition~\ref{prop:correction}, and $d=[K:\Q]$, then
$\OK$ has a definition with $d(N+1)$ total coordinates and
$d$ equations of degree at most
\[
 D'=|B|+2\max(D,|A|+1).
\]
It also has a single-equation definition in the same coordinates
of degree at most $dD'$.
\end{theorem}
\begin{proof}
Northcott's theorem makes $A$ finite. By hypothesis $B$ is finite,
and $A\cap B$ is empty because every rational integer is integral.
The corrected set
\[
 S_0=(S\setminus A)\cup B
\]
is Diophantine by Proposition~\ref{prop:correction} and satisfies
$\Z\subseteq S_0\subseteq\OK$.
Choose an integral basis $\omega_1,\ldots,\omega_d$ of $\OK$.
Then
\[
                    \OK=\sum_{i=1}^d\omega_iS_0.
\]
The inclusion from right to left uses integrality; the reverse
inclusion uses the integral-basis expansion with integer
coefficients, all of which belong to $S_0$.

For input $X$, retain $s_1,\ldots,s_{d-1}$ and substitute
\[
 s_d=\frac{X-\sum_{i<d}\omega_i s_i}{\omega_d}
\]
into the last of $d$ independent copies of the defining
polynomial for $S_0$. The substitutions are linear, and the
number of coordinates is
\[
 1+(d-1)+dN=d(N+1).
\]
This proves the stated system bound.
To combine its residuals $R_1,\ldots,R_d$ into one equation,
choose a degree-$d$ field extension $F/K$ with basis
$e_1,\ldots,e_d$ and use
\[
                  N_{F/K}\left(\sum_i e_iR_i\right)=0.
\]
A field element has norm zero exactly when it is zero, and
basis independence then forces every $R_i=0$. The norm is a
homogeneous polynomial of degree $d$ in the residuals, giving
the degree bound. Such an extension exists, for example by an
Eisenstein polynomial at a finite prime of $K$; for $d=1$ use
$F=K$.
\end{proof}

The theorem permits $S_0$ to contain nonrational algebraic
integers. Thus neither existential definability of $\Q$ in $K$
nor an exact index-equality conclusion is needed for this
criterion. The substantive remaining requirements are the
Diophantine definition of $S$, integer completeness outside a
finite set, and the height bound on its nonintegral inputs.

\begin{corollary}\label{cor:rationalintegers}
Suppose instead that $\Z\setminus S$ is finite and that every
$x\in S\setminus\Z$ has bounded height. Then $\Z$ is
Diophantine in $K$. With
$A=S\setminus\Z$ and $B=\Z\setminus S$, the polynomial and
bounds of Proposition~\ref{prop:correction} define $\Z$ itself.
\end{corollary}
\begin{proof}
Northcott's theorem makes $A$ finite, and
$(S\setminus A)\cup B=\Z$.
\end{proof}

\section{Use in height arguments}

If a proposed detector is complete for all nonzero integers but
its soundness argument only bounds the heights of false inputs,
Corollary~\ref{cor:rationalintegers} already gives the required
existence conclusion, with the possible omitted input zero
included in $B$. One need not force a contradiction for every
false input by introducing a further fixed parameter.
For a direct integrality detector, Theorem~\ref{thm:integers}
requires the height bound only on nonintegral inputs.

This observation does not assert that an available detector
satisfies these hypotheses over every number field. In
particular, estimates for elliptic point heights do not by
themselves prove local detector completeness or the existence
of rational witnesses in the required field. The finite
exception sets may also be ineffective even when their
finiteness is known. The assertions above concern existence of
a formula unless those sets are separately identified.

\begin{thebibliography}{9}
\bibitem{HS}
M. Hindry and J. H. Silverman,
\emph{Diophantine Geometry: An Introduction},
Graduate Texts in Mathematics 201, Springer, 2000,
Theorem B.2.3.
\bibitem{FieldDefinitions}
\emph{Existential Definitions of Rings of Integers and Subfields},
project manuscript, 9 November 2026,
\nolinkurl{papers/number_fields/number_field_definability_20261109_v12.tex}.
\end{thebibliography}
\end{document}
